\documentclass[12pt]{article}
\usepackage[utf8]{inputenc}
\usepackage[T1]{fontenc}
\usepackage{hyperref}
\title{Yauma Language Resources for K-12 Curriculum Development}
\author{Compiled by Claude Sonnet 5.5 (claude-sonnet-5-5) for Andr\'as Kornai}
\date{}
\begin{document}
\maketitle

\section{Introduction}
Yauma (ISO 639-3: yax; Glottocode yaum1238) is a Bantu language spoken in Angola and Zambia, mainly around the Kwando (Cuando) River. The sources found disagree on its status. One treats it as a variety of Mbunda (Guthrie K.10), and Ethnologue is reported to place it in the Ngangela subgroup. Speaker counts range from about 22,200 to 65,000. Searches found no Yauma-specific educational or reference material.

\section{Primary Educational Resources}
No primers, textbooks or school programmes specific to Yauma were found. The reference description says it is not known to be taught in schools. Materials in the closely related Mbunda may be relevant, but this search did not examine them.

\section{Linguistic and Reference Materials}
\begin{itemize}
\item Wikipedia entry for ISO 639:yax, which redirects to Mbunda language: \url{https://wikipedia.com/wiki/ISO_639:yax}
\item Endangered Languages Project entry: \url{https://endangeredlanguages.com/lang/5082}
\item Verbix endangered-languages page on Yauma: \url{https://docs.verbix.com/EndangeredLanguages/Yauma}
\end{itemize}

\section{Religious and Community Resources}
\begin{itemize}
\item Scripture Earth listing, which reports only portions of Scripture (the Gospel of John): \url{https://www.scriptureearth.org/00eng.php?iso=yax}
\item Joshua Project language profile: \url{https://www.joshuaproject.net/index.php/languages/yax}
\end{itemize}

\section{Digital and Audio Resources}
No audio, video or software resources specific to Yauma were found.

\section{Curriculum Development Context}
Yauma may be a dialect of Mbunda, so a curriculum may need to build on Mbunda materials and standards, subject to what the Yauma community prefers. The Gospel of John is the only text found. Developers would need to establish the orthography and write new materials with speakers in Angola and Zambia. Very little was found in these searches.

\end{document}